\section{Measurement discrimination, block resources, and learning}
\label{sec:matrixcomparison76}
\label{sec:blockcomparison78}
\label{sec:confidencecomparison80}

\subsection{Finite-use discrimination and horizontal geometry}
The measurement interface determines which discrimination results
apply.  Here the ordered classical outcome and an arbitrary external
reference are retained, while the device returns no postmeasurement
quantum system.  For projective measurements, Pucha\l a and
collaborators \cite[Corollary~1 and Theorem~2]{PPKKO72} give an
exact multiple-use formula with ancillary access, used in the
observable-readout theorem.  Fiur\'a\v sek--Mi\v cuda \cite{FM75}
study a restricted ancillary interface, with a different comparison
of adaptive and parallel probes.

For general POVMs, Krawiec--Pawela--Pucha\l a
\cite[Section~4]{KPP76} construct a pair of nine-outcome qutrit
rank-one measurements admitting perfect adaptive discrimination
with two calls and no perfect discrimination by any finite parallel
experiment.  Datta--Biswas--Saha--Augusiak
\cite[Theorems~3--5]{DBSA76} study entanglement-assisted single-use
discrimination, including perfect constructions and a qubit
entanglement advantage; their qubit construction has three outcomes.
These results concern retained-reference discrimination for broader
outcome alphabets.  For ordered binary effects,
\eqref{eq:matrixadaptivity76} bounds the adaptive advantage uniformly
in the horizon for each fixed input dimension.

The differential argument in Lemma~\ref{lem:matrixpath76} uses the
horizontal-gauge and channel-extension method
\cite{FI76,DKG67,DM76,KGAD76}.  Yuan--Fung give the parallel path
estimate \cite[Section~III, equation~(26), PDF version~3]{YF76},
and Sieniawski--Demkowicz-Dobrza\'nski integrate adaptive bounds
\cite[Section~IV.A]{SD76}.  The binary Sylvester choice produces
the regularized variance $M(I-M)$, whose path bound can be
integrated explicitly and matched by compressed nonadaptive tests.
Inverse-Sylvester expressions in Bures geometry
\cite{Dittmann76} and spectral matrix integration
\cite{SZGeom76} provide further antecedents.  Here the form acts on
the effect displacement $H$ with variance $V=M(I-M)$; the Bures
pullback under $M\mapsto V$ would instead act on the differential
$H-MH-HM$.  Theorem~\ref{thm:matrixcover76} combines the resulting
geometry with uniform operational-ball estimates and the nested
endpoint integral to obtain the binary-effect covering law.

\subsection{One-use tomography and acquisition resources}
Mele--Bittel \cite[Theorem~I.1]{MB67} give channel tomography with
$O(d_{\rm in}d_{\rm out}k\epsilon^{-2})$ calls at fixed confidence.
Their binary specialization
\cite[Lemmas~III.7--III.8 and Corollary~III.9]{MB67} returns a
legal effect with operator error at most $\epsilon$, using
\[
 M=1024d^2\epsilon^{-2}
   +O(d\epsilon^{-2}\log(1/\eta)),
 \qquad 4e^{-4d^2}<\eta<1,\quad 0<\epsilon<1.
\]
It prepares independent maximally entangled input--reference pairs
and processes the Choi copies collectively
\cite[Algorithm~I.2, Figure~1 and Section~III.2]{MB67}.  It
therefore belongs to the fresh-block class with $b=1$, which allows
coherent purification and collective output tomography.  The
source calls all invocations one parallel block; its inputs from
distinct invocations nevertheless remain independent.  Its
amplification \cite[Remark~III.4, equation~(165)]{MB67} gives
$O(d^2\epsilon^{-2}\log(1/\eta))$ calls for all
$0<\eta\le1/8$ with the same acquisitions.  The sharper confidence
extension derived below uses the dimension-dependent tail in the
printed Corollary~III.9.

Zambrano--Ramos-Calderer--Kueng
\cite[Definition~1 and Theorem~2]{ZRK76} use worst-input one-use
outcome total variation.  With $L$ outcomes, their global-$2$-design
protocol has sufficient budget
\[
 M\ge \frac{8(d^2+\epsilon(d^2+1)/6)}{\epsilon^2}
       \log\frac{2^{L+1}9^{2d}}{\eta}.
\]
It probes known inputs separately and returns a legal POVM by least
squares and projection.  Thus it also has $b=1$, without a retained
reference.  When $L=2$, its loss is $\|E-F\|_{\rm op}$, whereas
$d_1(E,F)=2\|E-F\|_{\rm op}$ in our normalization.  Their
Theorem~9 gives a constant-confidence lower bound
$\Omega((d^3+d^2L)\epsilon^{-2})$ for nonadaptive single-copy
measurements on known input states.

\subsection{Interior learning under the future-use loss}
Theorem~\ref{thm:interiorlearning78} applies the Mele--Bittel
binary estimator to $I/4\preceq E\preceq3I/4$ with
$\epsilon=\delta/(8\sqrt N)$ and $\eta=1/8$.  The horizontal
path estimate gives
\[
 d_N(E,F)\le4\sqrt N\|E-F\|_{\rm op}
 \qquad\text{when }\|E-F\|_{\rm op}\le1/8.
\]
Consequently $O(d^2N\delta^{-2})$ independent one-call probes
suffice at this confidence.  Theorem~\ref{thm:dimensionlower78}
gives the matching lower bound under arbitrary coherent adaptive
training.  This identifies the joint dimension, horizon, and
accuracy order on the stated interior.

For all integers $d,N\ge1$, $0<\delta\le2^{-13}$ and
$0<\eta\le1/8$, Theorem~\ref{thm:interiorconfidence80} gives
\[
 M^*_{\rm int}(d,N,\delta,\eta)
 =\Theta\!\left(N\delta^{-2}
             [d^2+d\log(1/\eta)]\right)
\]
with absolute constants.  The upper bound has $b=1$ and the lower
bound permits unrestricted coherent adaptive training.
Corollary~\ref{cor:operatorconfidence80} gives the corresponding
binary operator-learning order
$\Theta(\epsilon^{-2}[d^2+d\log(1/\eta)])$ for
$0<\epsilon\le2^{-14}$.

The upper bound uses the Mele--Bittel estimator with base failure
$q_d=2^{-3d}$, which satisfies $4e^{-4d^2}<q_d$ in their
Corollary~III.9.  Run it at a fixed fraction of the desired accuracy,
and repeat independently a number of times equal to the smallest
positive odd integer at least $2\log_2(1/\eta)/d$.  Metric-majority
selection and the binomial tail estimate then give
$O(\epsilon^{-2}[d^2+d\log(1/\eta)])$ calls for every
$0<\eta\le1/8$.  Lemma~\ref{lem:binaryconfidenceupper80}
records this confidence extension of the existing estimator.

For the confidence converse, Mele--Bittel's classicalisation of
diagonal multiple-coin channels
\cite[Lemma~IV.19 and Corollary~IV.20]{MB67} is a direct
antecedent.  Their corollary gives
$\Omega(d\epsilon^{-2}\log(d/\eta))$ for small operator error
and $0<\eta<1/4$.  Lemma~\ref{lem:diagonalconfidence80} gives a
null-versus-coordinate argument under the present coherent adaptive
interface and future-use loss.  Revealing the sampled diagonal
index isolates the hypothesis-dependent Bernoulli coin after a
common basis-dephasing map.  The classical--quantum relative-entropy
calculation permits different current quantum inputs under the
two hypotheses.  Summation of the testing inequalities gives
$\Omega(dN\delta^{-2}\log(1/\eta))$; its maximum with
Theorem~\ref{thm:dimensionlower78} is the asserted sum up to an
absolute factor.

Channel-learning converses have different domains.
Oufkir--Girardi \cite{OG67} and Chen--Zhang--Yu
\cite[Theorems~1.1 and~4.2]{CZY78} treat fully coherent channel
learning in one-use diamond loss.  The latter gives
$\Omega(rd_{\rm in}d_{\rm out}\epsilon^{-2})$ when
$rd_{\rm out}\ge2d_{\rm in}$ at constant confidence.  Quantum
output dimension two still allows nonclassical outputs, so that
general-channel theorem does not specialize directly to a
binary-qc converse.  For binary measurements, Mele--Bittel give
their projector packing \cite[Proposition~IV.17]{MB67} and the
diagonal-coin bound above.  The near-fair binary information
estimate in Lemma~\ref{lem:weakmeasurementinfo78} supplies the
self-contained growing-dimension converse used here.

\subsection{The closed body and independent acquisition blocks}
On the complete effect body, the hybrid inequality
\[
 d_N(E,F)\le2N\|E-F\|_{\rm op}
\]
converts both one-use tomography guarantees to sufficient budgets
$O_d(N^2\delta^{-2}\log(1/\eta))$.  Theorem~\ref{thm:blocklearning78}
and Corollary~\ref{cor:onecallseparation78} give the corresponding
acquisition obstruction: for fixed $d\ge2$, every fresh $b=1$
procedure learning the entire ordered binary effect body requires
\[
 M\ge cN^2\delta^{-2}\log(1/\eta).
\]
The converse allows classical adaptive choices, stored references,
their collective processing, and public bounded stopping.  It
therefore applies to every reconstruction from the acquisitions of
both cited tomography protocols.  The projective boundary forces
this full-body horizon order, while the interior calculation above
permits the linear order using the same acquisition resource.

More generally, for $d\ge2$, $1\le b\le N$,
$0<\delta\le\delta_*$ and $0<\eta\le1/8$,
Theorem~\ref{thm:blocklearning78} proves
\[
 c\frac{N^2}{b}\delta^{-2}\log(1/\eta)
 \le M_b^*(d,N,\delta,\eta)
 \le Cd^4\frac{N^2}{b}\delta^{-2}\log(d/\eta).
\]
Here $M_b^*$ is the worst-record budget for a common classical
estimate in the fresh-block model of
Section~\ref{sec:blocklearning78}.  Each block may use a reference
and adaptive quantum controls.  Stored outputs can be processed
jointly, but only the classical record determines the next fresh
probe; earlier quantum registers cannot feed its unfinished queries.
At $b=N$ this gives the fixed-dimensional linear-horizon order,
and at $b=1$ the quadratic-horizon order.  For $d=1$, the cost is
$\Theta(N\delta^{-2}\log(1/\eta))$ independently of $b$.

The all-confidence operator estimator gives a further full-body
upper bound.  Use Lemma~\ref{lem:binaryconfidenceupper80} at
$\epsilon=\delta/(2N)$, or use the block construction, selecting
the smaller public budget before acquisition.  This yields
\eqref{eq:fullbodyconfidenceupper80}:
\[
 M_b^*(d,N,\delta,\eta)
 \le C\frac{N^2}{\delta^2}
       \min\!\left\{d^2+d\log\frac1\eta,
                    \frac{d^4}{b}\log\frac d\eta\right\},
\]
for $d\ge2$, $1\le b\le N$,
$0<\delta\le\min\{2^{-13},\delta_*\}$ and
$0<\eta\le1/8$.  The lower and upper dimension dependences for
the complete body remain unmatched when the parameters grow jointly.

The block scale has established metrological antecedents.
Hyllus et al.\ \cite[Observation~1, equation~(9)]{Hyllus78} and
T\'oth \cite[Observation~3, equation~(6)]{Toth78} prove, for
$b$-producible qubit probes under linear phase encoding,
\[
 F_Q\le\lfloor M/b\rfloor b^2+
              (M-b\lfloor M/b\rfloor)^2\le Mb.
\]
Our block parameter also permits sequential controls within a fresh
block and concerns query acquisition rather than multipartite
entanglement of the final state.  The finite-pair fidelity argument
gives the confidence factor without an unbiasedness assumption.
Salek--Hayashi--Winter \cite[Section~II.2 and Section~V.2]{SHW68}
formulate classical feed-forward with arbitrary output quantum
memory and no input quantum memory, and study asymptotic
discrimination exponents.  Wilde et al.\
\cite[Proposition~21 and Appendix~B]{WBHK72} give adaptive
divergence converses and fidelity-to-error inequalities.
Section~\ref{sec:blocklearning78} establishes the finite-budget
conditional fidelity estimate with variable block lengths and
stopping.  Its upper bound blocks the future tester and applies the
common learner at the shorter horizon.

Restrictions on jointly measuring already supplied state copies
are a separate resource \cite{CLL78,NZ78}.  The fresh-query
restriction used here permits an unrestricted final joint
measurement.  Mele--Bittel's blockwise Choi processing
\cite[Section~I.5]{MB67} concerns the former resource.

The full-body common learner uses multiscale phase recovery
\cite{HB76,KLY76}, with a proved safe-angle induction, bias
cancellation, unknown-noise block selection, and conditional
confidence allocation.  Its matrix extension calibrates both
support endpoints and controls variance from data-selected
compressions before assembling a legal estimate.  These operations
provide one procedure at rank changes and repeated eigenvalues;
the pair-dependent witnesses of the metric proof serve only to
establish operational separation.

\subsection{Exact descriptions and finite collective synthesis}
Theorem~\ref{thm:matrixcodec77} gives an exact rational selection
rule with optimal fixed-dimensional payload on the closed body
in its small-error range.  Theorem~\ref{thm:interiorentropy79}
has explicit dimension-uniform constants on $[I/4,3I/4]$.
Its norm-ball volume argument and the Fano--Kraft proof of
Theorem~\ref{thm:prefixrate79} use standard geometric and
information-theoretic mechanisms to obtain the future-loss
description laws.  Theorem~\ref{thm:interiorlearnedcode79}
coarsens the Mele--Bittel estimator to a finite output and the exact
public dictionary without further device calls; its confidence
extension is Corollary~\ref{cor:interiorconfidencecode80}.

Computing the finite readout is a further question: measurable
finite coarsening alone does not give computable measurement
entries.  In Section~\ref{sec:effectivelearning80},
Lemma~\ref{lem:algebraicreadout80} expresses the risk on
independent Choi records by a first-order formula over the reals
with rational coefficients.  It includes positivity and
completeness of the conditional POVMs, the fixed dictionary, and
quantification over every real target in the interior.  Real
quantifier elimination decides feasibility; real-closed-field
transfer supplies an algebraic point, found by exact enumeration
and sign tests \cite[Theorems~2.77 and~2.80 and Chapter~14]{BPR80}.
An increasing search of the sample count terminates at the order
supplied by the statistical existence bound.  The resulting
Theorem~\ref{thm:effectiveinterior80} gives an algebraic readout
and finite trusted controls with
\[
 M=O(N\delta^{-2}[d^2+d\log(1/\eta)]),\qquad
 B=\tfrac12d^2\log_2N+d^2\log_2(1/\delta)+O(d^2).
\]

Section~\ref{sec:finiterisk81} gives finite rational certificates
for the same uniform risk.  If a fixed readout has failure at most
$\alpha$ on an operator-norm $r$-net for good-label radius $a$,
Theorem~\ref{thm:finitecertificate81} gives
\[
 \sup_{E\in\mathfrak I_d}
   \mathbb P_E\{\|E-C_J\|_{\rm op}>a+r\}\le\alpha+mr.
\]
The analytic input is
$\|\Gamma_m(E)-\Gamma_m(G)\|_1\le2m\|E-G\|_{\rm op}$.
For each nearby grid point $G$, one transfers its fixed set of good
labels between the two experiments; no continuity of label
selection is assumed.  With
$a=\delta/(8\lceil\sqrt N\rceil)$ and
$r\le\min\{\delta/(32\lceil\sqrt N\rceil),\eta/(16m)\}$,
exact grid inequalities at failure $\eta/4$ certify operational
error $5\delta/8$ with failure at most $5\eta/16$ for all real
interior effects.

Proposition~\ref{prop:rationalreadout81} mixes an arbitrary
$L$-outcome readout with the uniform outcome and rounds its first
$L-1$ effects at prescribed denominator
$Q=\lceil64L^2d^m/\eta\rceil$.  The last effect is their
complement, giving exact completeness; the mixing margin preserves
positivity.  The change of every output-index law is bounded by
$3\eta/32$.  A finite enumeration therefore contains a certified
rational readout whenever the statistical construction supplies the
required slack.  Exhausting this enumeration at each dyadic call
count gives Theorem~\ref{thm:rationallearner81} at the same joint
call and payload orders.  Only existence of the tomography readout
is used: its entries and continuous-outcome integrals are not
inputs to either synthesis.

Finite-control approximation charges the complete preparation and
classically conditioned readout instruments.  The pathwise total
unhalved error allowance $\eta$ changes index probabilities by at
most $\eta/2$, so the rational construction has actual failure at
most $13\eta/16<\eta$.  Acquisition remains fresh and one-call,
and the decoder returns the unchanged legal rational centre.
Neither construction requires coherent access to the device's
classical outputs or environment.  Rational verification replaces
real quantifier elimination in the latter construction, while both
separate device calls and payload from preprocessing, storage, and
known-control work.

Finite and efficient state tomography have substantial existing
constructions.  Pelecanos--Spilecki--Tang--Wright
\cite[Theorems~1.1--1.3, Section~3.3 and Figures~8--9]{PSTW80}
give an efficient purification reduction and efficient pure- and
mixed-state tomography; their Theorem~4.2 supplies discrete
approximate-design circuits used by the pure-state protocol.
Preserving the covariance and binary operator-error analysis of the
imported channel estimator under a different tomography
implementation requires an additional argument.  The constructions
here instead select a feasible collective readout with a uniform
future-loss guarantee and the exact public dictionary.  They prove
finite computability and, for the rational construction, finite
verifiability of the risk conditions.  Their exhaustive search,
dictionary reconstruction, workspace, and known-gate costs have no
efficiency bound.

\begingroup
\small
\setlength{\tabcolsep}{4pt}
\renewcommand{\arraystretch}{1.17}
\begin{longtable}{@{}>{\raggedright\arraybackslash}p{0.25\textwidth}>{\raggedright\arraybackslash}p{0.35\textwidth}>{\raggedright\arraybackslash}p{0.34\textwidth}@{}}
\toprule
Result and acquisition & Loss and parameter range & Training-call order\\
\midrule
\endfirsthead
\toprule
Result and acquisition & Loss and parameter range & Training-call order\\
\midrule
\endhead
Mele--Bittel with Lemma~\ref{lem:binaryconfidenceupper80};
independent Choi probes, $b=1$ &
Entire binary body; operator error $\epsilon$;
$0<\eta\le1/8$ &
$O(\epsilon^{-2}[d^2+d\log(1/\eta)])$;
source estimator with a confidence extension\\
Zambrano et al.; individual known probes, $b=1$ &
Binary global-design protocol; operator error $\epsilon$;
$0<\eta<1$ &
$O(d^2\epsilon^{-2}[d+\log(1/\eta)])$\\
Theorem~\ref{thm:interiorconfidence80}; independent Choi probes,
$b=1$ &
$I/4\preceq E\preceq3I/4$; future error $\delta$;
$d,N\ge1$, $0<\eta\le1/8$ &
$\Theta(N\delta^{-2}[d^2+d\log(1/\eta)])$;
coherent adaptive lower bound\\
Theorem~\ref{thm:effectiveinterior80}; synthesized algebraic
readout, $b=1$ &
Same interior and future loss; rational $\delta,\eta$ &
$O(N\delta^{-2}[d^2+d\log(1/\eta)])$;
finite trusted-control specification\\
Theorem~\ref{thm:rationallearner81}; certified rational readout,
$b=1$ &
Same interior and future loss; rational $\delta,\eta$ &
$O(N\delta^{-2}[d^2+d\log(1/\eta)])$;
finite rational risk certificate\\
Theorem~\ref{thm:blocklearning78}; fresh blocks of size $\le b$ &
Entire binary body; future error $\delta$;
$d\ge2$, $1\le b\le N$, $0<\eta\le1/8$ &
$\Theta_d((N^2/b)\delta^{-2}\log(1/\eta))$;
explicit upper bound in
\eqref{eq:fullbodyconfidenceupper80}\\
\bottomrule
\end{longtable}
\endgroup
The future-loss rows use the absolute small-error ranges stated in
their theorems.  The sharp joint dimension law applies to the
displayed interior, while the closed-body law is optimal in
horizon, accuracy, and confidence for fixed dimension.  The binary
consuming interface is common to these statements; multiple-outcome
measurements and residual quantum outputs require a different
analysis.
